\chapter{Philosophy}


\section{Podcast -- Charlotte Luyckx, philosophe et chargée de cours invitée à l'UCLouvain}

\subsubsection{Que penses-tu du mot « complexe » ? \textnormal{\small(3:36)}}

Charlotte Luyckx utilise ce mot dans le sens qu'Edgar Morin lui donne, en revenant à sa racine étymologique \textit{complexus}, qui signifie « tisser ensemble ». Dans les questions énergétiques et environnementales, il s'agit effectivement de tisser ensemble des éléments qui ont souvent été pensés séparément, en silo, voire en opposition. La crise écologique invite à les penser conjointement et à les réconcilier.

\subsubsection{Peux-tu m'en dire un peu plus sur la théorie du doughnut ? \textnormal{\small(4:42)}}

La théorie du doughnut, développée par Kate Raworth, propose un modèle visuel qui réconcilie deux champs souvent pensés séparément : la justice sociale et l'écologie. Le modèle identifie d'une part des \textbf{plafonds environnementaux} à ne pas dépasser -- fondés sur les neuf limites planétaires identifiées par le collectif de scientifiques autour de Rockström (climat, biodiversité, éléments chimiques dans l'atmosphère, utilisation des terres, etc.) -- et d'autre part des \textbf{planchers sociaux} en dessous desquels on ne peut descendre (accès à l'éducation, au logement, à la nourriture, à un revenu). L'enjeu est de se maintenir dans la zone intermédiaire, entre plancher et plafond. Or, à l'heure actuelle, on est dans le rouge tant au niveau des plafonds que des planchers dans de nombreux domaines, même si la situation varie selon les zones géographiques. Le modèle fonctionne comme une boussole ou un tableau de bord : il ne résout rien en lui-même, mais permet de visualiser deux problèmes conjoints qui, traditionnellement, ne sont pas pensés ensemble.

\subsubsection{Peux-tu me parler un peu de l'éco-philosophie ? \textnormal{\small(8:01)}}

L'éco-philosophie est la rencontre des questions environnementales et écologiques avec la philosophie. On pourrait la définir comme l'ensemble des enjeux philosophiques de la crise écologique. On trouve ce concept notamment chez le philosophe Arne Næss. L'éco-philosophie permet de penser l'écologie non seulement comme un problème technique ou politique, mais comme une question soulevant des interrogations philosophiques fondamentales. La relation est à double sens : d'un côté, les enjeux philosophiques offrent un recul et une profondeur sur les questions écologiques ; de l'autre, ces questions interpellent la manière même de pratiquer la philosophie aujourd'hui et peuvent renouveler, voire rafraîchir certains concepts, en revisitant la tradition philosophique à l'aune de la crise écologique -- à l'aune de la question de l'\textit{oïkos}, de cette maison que nous habitons et qui n'a pas toujours été placée à sa juste place dans notre compréhension de nous-mêmes et du monde.

\subsubsection{Quelles sont les différentes strates ? \textnormal{\small(12:22)}}

Le modèle des strates vise non pas à hiérarchiser les dimensions de la crise, mais à identifier chacune dans sa pertinence propre, en soulignant leur complémentarité. Les strates les plus visibles (plus en surface) sont les mieux représentées dans les médias et les débats publics ; les strates plus profondes nécessitent davantage d'effort pour être mises en lumière. Voici les strates, de la plus visible à la plus profonde :

\begin{enumerate}
    \item \textbf{Les enjeux techniques et technologiques} : c'est la réponse la plus immédiate et la plus médiatisée à la crise énergétique et écologique. Elle est pertinente mais insuffisante, notamment en raison des effets rebonds et de l'énergie grise.
    \item \textbf{Les comportements individuels et collectifs} : sortir du consumérisme, changer son alimentation, ses modes de déplacement. Cette strate est essentielle, mais limitée : selon l'étude de Carbone 4, environ 60\,\% de l'empreinte carbone individuelle est structurelle.
    \item \textbf{Les enjeux structurels économiques et politiques} : changer les modèles économiques, les indicateurs de croissance, voire les axes de la pensée politique (comme le propose Bruno Latour dans \textit{Où atterrir?}).
    \item \textbf{L'éco-philosophie} : redéfinir des concepts comme le progrès ou le développement ; interroger la place de l'humain dans le cosmos ; aborder des questions d'éthique environnementale et d'ontologie.
    \item \textbf{L'éco-spiritualité et l'éco-théologie} : questionner le système de croyances qui sous-tend notre culture, les sources de sens et d'engagement face à la crise, et la dimension intérieure de la transition.
\end{enumerate}

\subsubsection{Si on continue à creuser après les deux premières strates, que trouve-t-on ? \textnormal{\small(18:00)}}

En creusant au-delà des strates technologique et comportementale, on arrive aux \textbf{enjeux structurels}, principalement économiques et politiques, qui visent un changement des modèles économiques, des indicateurs de croissance et des axes de la pensée politique. Bruno Latour, dans \textit{Où atterrir?}, invite par exemple à repenser l'axe même de l'engagement politique, en passant du modèle moderne à une orientation vers ce qu'il appelle « le terrestre ».

En creusant encore, on atteint le \textbf{système de croyances} qui a contribué à forger nos institutions. Des questions éco-philosophiques émergent alors : redéfinition du progrès et du développement, place de l'humain dans le cosmos, choix du vocabulaire pour désigner notre être au monde, questions d'éthique environnementale et de justice, et enfin questions ontologiques ou métaphysiques sur la nature même du réel.

La strate la plus profonde est la strate \textbf{éco-spirituelle ou éco-théologique}, qui interroge les racines religieuses et spirituelles de notre rapport à la nature. L'historien médiéviste Lynn White incrimine notamment la théologie transcendentaliste du christianisme médiéval -- qui a externalisé Dieu par rapport au monde créé, et par miroir l'humain par rapport à la nature -- comme matrice conceptuelle ayant préparé l'avènement de la modernité et contribué à la marchandisation du monde.

\subsubsection{Peux-tu nous parler de la strate comportementale et de l'effet colibri ? \textnormal{\small(24:54)}}

La strate comportementale mise sur les changements individuels et collectifs : sobriété, alimentation, modes de déplacement, zéro déchet. Elle est illustrée par la métaphore du colibri de Pierre Rabhi : face à l'incendie de la forêt, le petit colibri apporte ses trois gouttes d'eau en disant « je fais ma part ». Cette approche représente une forme d'engagement moral, une éthique non strictement conséquentialiste, qui consiste à agir conformément à sa conscience indépendamment de l'impact sociétal mesurable.

Cette strate est fondamentale : elle crée de nouvelles pratiques sociales et des alternatives aux modèles dominants. Mais elle est insuffisante comme seule réponse, pour au moins deux raisons. Premièrement, une part importante de l'empreinte carbone (environ 60\,\% selon Carbone 4) est structurelle et échappe aux choix individuels. Deuxièmement, la métaphore du colibri comporte une dimension individualiste : aller seul mettre de l'eau sur le feu, sans recréer du collectif et du politique. Il faudrait que le colibri réunisse les autres animaux pour agir ensemble. Par ailleurs, cette logique peut faire peser sur les individus des responsabilités qui relèvent en réalité du niveau structurel, et peut culpabiliser des personnes dont les moyens ne leur permettent pas de choisir une consommation plus écologique.

\subsubsection{C'est intéressant de voir les différentes portes d'entrée vers un mode de vie plus respectueux de notre planète. \textnormal{\small(32:56)}}

Charlotte Luyckx souligne qu'il ne s'agit pas de choisir entre les différentes formes d'engagement -- réduire sa consommation de viande, limiter l'usage de la voiture, pratiquer le zéro déchet -- mais de les additionner plutôt que de les mettre en opposition. Des recherches en sociologie (notamment celles d'Émeline de Bouver sur la simplicité volontaire) montrent d'ailleurs que les personnes engagées dans un changement de comportement ne sont généralement pas enfermées dans cette seule strate : elles participent aussi aux manifestations pour le climat et tentent de créer du collectif dans leur quartier. Les différentes portes d'entrée sont donc complémentaires, et il est important d'objectiver les impacts relatifs de chaque geste -- le régime végétarien ayant par exemple un impact bien plus significatif que de nombreux autres écogestes -- pour mieux orienter les efforts sans culpabiliser.

\subsubsection{Qui a dit la phrase suivante ? « Parce qu'elle est la clé du progrès environnemental, parce qu'elle fournit les ressources permettant d'investir dans les technologies propres, la croissance est la solution, non le problème. » \textnormal{\small(34:53)}}

Cette phrase a été dite par \textbf{George W. Bush}. Elle illustre un mythe extrêmement tenace, partagé par de nombreux dirigeants et économistes : celui de la croissance comme condition de possibilité des démocraties et du progrès environnemental lui-même.

\subsubsection{Que penses-tu du mythe de la croissance infinie ? \textnormal{\small(36:01)}}

Charlotte Luyckx rappelle que depuis au moins cinquante ans, des théoriciens ont montré clairement qu'une croissance infinie dans un monde fini est une impossibilité. Des économistes comme Nicholas Georgescu-Roegen (bioéconomie) ont croisé données environnementales et modèles économiques pour démontrer cette contradiction. Ivan Illich, dès les années 1970-1980, insistait sur la nécessité d'établir des seuils : des études montrent qu'au-delà d'un certain niveau de confort matériel, la corrélation entre croissance du PIB et bien-être disparaît, voire s'inverse. L'enjeu aujourd'hui est d'intégrer dans notre culture cette notion de limite et de seuil, que notre culture du progrès indéfini a perdue. La croissance infinie relève davantage du \textbf{mythe} que de la connaissance scientifique : c'est une croyance nourrie par des éléments philosophiques de la modernité qui ont alimenté cette fuite en avant vers le progrès matériel.

\subsubsection{Tu utilises souvent Banksy pour illustrer les questions sur le PIB. Pourquoi cet artiste ? \textnormal{\small(40:03)}}

Charlotte Luyckx apprécie Banksy pour plusieurs raisons. D'abord esthétiquement, l'œuvre lui parle, notamment grâce au mystère entourant l'artiste. Ensuite, comme manière à la fois humoristique et poétique de critiquer le capitalisme, elle trouve la démarche remarquable. Banksy met également en lumière le paradoxe inhérent à toute critique du capitalisme : on est soi-même inscrit dans le système qu'on critique, une tension que l'artiste rend visible avec acuité.

\subsubsection{C'est difficile de critiquer le capitalisme, vu qu'on en fait partie. \textnormal{\small(41:42)}}

Ce paradoxe peut être illustré par l'image du \textbf{baron de Münchhausen}, qui se sort d'un marécage en tirant sur ses propres cheveux. Comme l'ont montré des penseurs de la décroissance tels que Serge Latouche et Aminata Traoré, l'idéologie de la croissance et du progrès n'est pas seulement imposée de l'extérieur : elle a \textbf{colonisé nos imaginaires}. S'y opposer implique donc un travail sur soi-même. L'économiste Christian Arnsperger parle du « capitaliste intérieur » avec lequel il faut apprendre à composer. Derrière l'aspiration à la croissance illimitée, il y a peut-être une soif légitime d'infini et d'épanouissement -- mais il faudrait lui donner d'autres formes que le capitalisme : l'art, l'affinement des capacités perceptives, la contemplation du monde vivant.

\subsubsection{Tu dis que le développement durable pourrait être un oxymore. Peux-tu m'expliquer ? \textnormal{\small(44:34)}}

Selon Charlotte Luyckx -- en s'appuyant notamment sur Serge Latouche -- l'expression « développement durable » ou « croissance soutenable » est un oxymore, car le développement tel qu'il est conçu dans le paradigme économique moderne est intrinsèquement porteur d'un projet contraire aux exigences de réinsertion de nos activités dans les limites planétaires. Associé à l'illimitisme et à la croissance, il ne peut structurellement pas être durable. L'expression masque donc une contradiction fondamentale entre deux termes incompatibles dans leurs définitions courantes.

\subsubsection{Est-ce que le terme « post-croissance » est une manière politiquement correcte de dire « décroissance » ? \textnormal{\small(45:47)}}

Pas uniquement. L'économiste Isabelle Cassiers (UCLouvain) préfère le terme \textbf{post-croissance} pour des raisons à la fois rhétoriques et conceptuelles. Rhétoriquement, « décroissance » est peu vendable comme projet collectif et fonctionne facilement comme repoussoir. Conceptuellement, la décroissance est la photo négative de la croissance : elle manque les distinctions qualitatives nécessaires. Tout ne doit pas croître indéfiniment, mais tout ne doit pas non plus décroître indéfiniment. La post-croissance désigne plutôt \textbf{une sortie de la logique de la croissance} : évacuer la réponse automatique selon laquelle croissance = prospérité = bonheur, tout en se demandant ce qui doit croître, dans quel système de valeurs, et jusqu'à quel seuil.

\subsubsection{Quels seraient des exemples de bons indicateurs pour représenter la croissance ? \textnormal{\small(47:24)}}

Charlotte Luyckx renvoie aux travaux d'Isabelle Cassiers, qui a consacré une grande partie de sa carrière à ces indicateurs alternatifs. Elle cite notamment le modèle bhoutanais du \textbf{Bonheur Intérieur Brut (BIB)}, développé avec un panel d'experts internationaux dont Cassiers faisait partie, qui propose une palette d'indicateurs alternatifs visant à redéfinir collectivement ce qui constitue une vie réussie. Cette démarche soulève des questions fondamentales : qu'est-ce que la prospérité, le bonheur, une vie humaine épanouie ? Les réponses sont inexorablement plurielles, ce que masquait mal la réponse automatique par la croissance du PIB. La difficulté est alors de faire cohabiter cette pluralité des visions de la vie bonne avec une orientation politique commune -- une tension philosophique majeure.

\subsubsection{Les autres indicateurs que la croissance engendrent de nouvelles questions : qu'est-ce que la justice, le bien et le mal ? T'aventures-tu aussi dans les réponses ? \textnormal{\small(49:14)}}

Charlotte Luyckx ne prétend pas avoir répondu à toutes ces questions philosophiques. Elle les nomme et en identifie les enjeux, mais renvoie volontiers à des spécialistes comme Axel Gosseries (UCLouvain) pour les questions de justice environnementale. La question qui retient particulièrement son attention est celle du \textbf{rapport à la nature} : quelle représentation de la nature, et comment comprendre notre situation humaine dans un espace terrestre où nous cohabitons avec le non-humain et d'autres formes de vie ? Cette question constitue l'un des fils conducteurs de son travail éco-philosophique.

\subsubsection{Tu mentionnais qu'il y avait un « deep » et un « shallow ». Que trouve-t-on si on continue à creuser ? \textnormal{\small(55:31)}}

En continuant à creuser au-delà de l'éco-philosophie, on atteint la strate \textbf{éco-spirituelle ou éco-théologique}. Cette strate est plus large qu'une simple réflexion sur Dieu : l'éco-spiritualité englobe toutes les questions de sens qui émergent face à l'état du monde, les ressources intérieures permettant de s'engager, les finalités d'une nouvelle culture à construire. Elle inclut l'idée qu'une \textbf{transition intérieure} est nécessaire en plus de la transition extérieure -- « se changer soi pour changer le monde » (Gandhi). Cela passe par des exercices de reconnexion avec la nature, un changement d'échelle temporelle, et une réappropriation de notre rapport au vivant et à notre propre corps. On trouve aussi dans ce champ des relectures éco-théologiques des grandes traditions (christianisme, islam, bouddhisme, hindouisme), posant la dialectique entre immanence et transcendance dans leur rapport à la nature.

\subsubsection{Quel serait l'impact que pourrait avoir l'encyclique du pape François ? \textnormal{\small(59:37)}}

L'encyclique \textit{Laudato si'} (2015) représente, aux yeux de Charlotte Luyckx, l'un des textes de sensibilisation écologique à plus fort impact potentiel. Environ 25\,\% de la population mondiale est chrétienne, et une part significative reconnaît l'autorité du pape. À partir du moment où celui-ci prend une position courageuse et ambitieuse sur la crise écologique, le texte peut toucher un public bien plus large que les rapports du GIEC, réservés à une élite disposant des outils pour les lire. L'encyclique peut ainsi percoler dans des milieux qui ne sont pas directement connectés au débat scientifique, et avoir un impact politique non négligeable à l'échelle mondiale.

\subsubsection{Y a-t-il un élément en particulier que tu trouves intéressant dans l'encyclique ? \textnormal{\small(1:00:49)}}

Lue à travers la grille des strates, l'encyclique comporte, selon Charlotte Luyckx, plusieurs passages qui répondent directement à la critique de Lynn White adressée au christianisme. Ces passages affirment la \textbf{présence immanente de Dieu dans le monde créé}, ce qui constitue une position théologiquement audacieuse. Cela conduit à une vision que l'on pourrait qualifier de \textbf{panenthéiste} : Dieu n'est pas totalement extérieur à sa création (transcendantalisme pur), ni confondu avec elle (immanentisme pur), mais il l'\textit{habite}. Cette synthèse entre transcendance et immanence ouvre une voie vers une reconsidération de la sacralité du monde naturel au sein même de la tradition chrétienne.

\subsubsection{En quoi les modèles philosophiques peuvent-ils nous aider à prendre des actions ? \textnormal{\small(1:02:58)}}

L'idée sous-jacente au modèle des strates est qu'à chacune des strates correspond une forme d'engagement et d'action possible. Le « faire » a plusieurs visages : installer des panneaux photovoltaïques, bien sûr, mais aussi réfléchir à des indicateurs alternatifs au PIB, s'engager dans une redéfinition de notre rapport aux non-humains, prendre le temps de l'expérience perceptive de la nature, ou interroger ses propres aspirations et angoisses. Tout cela contribue à la mise en mouvement vers la transition. L'enjeu est donc de voir la \textbf{complémentarité} entre les différentes formes d'engagement plutôt que de les mettre en compétition -- éviter les guerres de chapelle entre approches technicistes et approches philosophiques, entre engagement individuel et action politique. Un changement de culture nécessite des transformations à tous les niveaux simultanément.

\subsubsection{Faut-il une vision intégrée ? \textnormal{\small(1:07:34)}}

Idéalement, oui -- mais sans prétendre que chacun peut tout faire. Charlotte Luyckx et la sociologue Émeline de Bouver ont co-écrit un article intitulé « L'écosystème de la transition », qui développe cette idée : chaque acteur peut s'inscrire dans une dimension de la crise (technologique, philosophique, sociale, etc.) et contribuer à un écosystème plus large, en cherchant des synergies plutôt qu'en entrant en compétition. L'\textbf{écologie intégrale} se définit ainsi comme une vision intégrative qui reconnaît la pluridimensionnalité de la crise et articule ses différentes dimensions, sans perdre d'énergie dans des oppositions stériles.

\subsubsection{Comment répondre à une urgence tout en laissant le doute ? \textnormal{\small(1:08:44)}}

Charlotte Luyckx distingue deux types de doute. Il y a le \textbf{doute scientifique légitime}, constitutif de la démarche épistémique, et le \textbf{doute pseudo-scientifique orchestré} par des intérêts privés -- les « marchands de doute » décrits par Naomi Oreskes, qui fabriquent médiatiquement une incertitude artificielle au service d'intérêts économiques (comme ce fut le cas avec le tabac ou le sucre). Il importe de mieux faire connaître la spécificité épistémique de la démarche scientifique -- sans tomber dans un positivisme naïf -- pour que des institutions comme le GIEC puissent avoir un poids politique fondé sur leur ancrage méthodologique. Quant au sentiment d'urgence face à l'ampleur des changements nécessaires et à la lenteur des transformations, Charlotte Luyckx l'assume et y répond par une posture proche de celle du colibri : faire ce qui est en son pouvoir, à toutes les strates, sans se laisser glisser vers un fatalisme paralysant.

\subsubsection{À l'inverse de l'éco-totalitarisme/fascisme, faudrait-il plus de transparence et de démocratie ? \textnormal{\small(1:16:18)}}

Charlotte Luyckx se positionne clairement en faveur d'une \textbf{démocratie écologique}, tout en reconnaissant ses limites intrinsèques. Elle rejette l'idée d'un gouvernement autoritaire, pour deux raisons principales : d'abord, la création culturelle profonde nécessaire à la transition ne peut pas être télécommandée -- elle exige une part irréductible de dynamiques bottom-up, d'innovation sociale et culturelle qu'une approche totalitaire empêcherait. Ensuite, l'étiquette « éco-fasciste » a souvent été utilisée abusivement comme repoussoir pour discréditer des courants de pensée écologique qui valorisent au contraire la diversité naturelle et culturelle. Elle s'appuie sur Dominique Bourg (\textit{Pour une démocratie écologique}) pour penser une démocratie rénovée, purgée de certains travers modernes : meilleure intégration des échelles de temps et d'espace, et plus grande transparence dans les processus de décision.

\subsubsection{Il faudrait trouver une manière d'intégrer les générations futures dans la démocratie ? \textnormal{\small(1:19:58)}}

Oui, c'est un des défis centraux d'une démocratie écologique. Les générations futures et le non-humain constituent de nouveaux acteurs politiques de l'Anthropocène qui ne peuvent pas s'exprimer par eux-mêmes. Il faudrait leur trouver des représentants ou tuteurs pour qu'ils puissent peser dans les débats. Le philosophe \textbf{Hans Jonas}, dans son \textit{Principe Responsabilité}, a posé la question du statut juridique et moral des êtres humains à naître. Une démocratie écologique devrait également intégrer le \textbf{passé} dans son rapport au temps long : les pays occidentaux émettent du CO\textsubscript{2} depuis la révolution industrielle, quand d'autres n'en émettent que depuis deux décennies. Les questions de justice climatique -- y compris les mécanismes de réparation et la dimension coloniale -- sont donc indissociables de cette réflexion sur la démocratie à venir.

\subsubsection{Qu'est-ce qui nous empêche de mettre les choses en œuvre ? \textnormal{\small(1:21:13)}}

Plusieurs obstacles sont identifiés. D'abord, la \textbf{non-perceptibilité} de certaines crises : si les canicules et inondations rendent le dérèglement climatique plus tangible, l'érosion de la biodiversité ou l'utilisation des terres restent abstraites pour le citoyen urbain lambda. Ensuite, la \textbf{profondeur de la crise culturelle} elle-même : elle ébranle notre système de croyances à un niveau fondamental, ce qui génère une résistance existentielle -- on est attaché à ce qu'on connaît, et changer de modèle implique de renoncer à des choses valorisées. La \textbf{colonisation de l'imaginaire} par l'idéologie de la croissance rend difficile d'envisager autre chose. Enfin, il y a des \textbf{intérêts privés} qui s'investissent activement contre le changement, parfois justifiés par des impératifs d'emploi ou d'approvisionnement. Charlotte Luyckx préfère une lecture bienveillante -- erreur ou manque d'information plutôt que mauvaise foi -- tout en reconnaissant que certains dirigeants semblent, comme le suggère Bruno Latour dans \textit{Où atterrir?}, avoir renoncé au projet du bien commun.

\subsubsection{Quel est le mauvais conseil que tu entends souvent ? \textnormal{\small(1:29:46)}}

Charlotte Luyckx nuance l'idée même de « mauvais conseil » : un conseil n'est pas mauvais en soi, mais peut l'être s'il est mal adapté au contexte ou à la personne. Le même conseil -- « engagez-vous davantage sur les questions environnementales » -- sera pertinent pour quelqu'un peu sensibilisé, et dramatique pour un militant en burn-out. Il s'agit donc plutôt d'adéquation au contexte que de mauvais conseils absolus.

\subsubsection{Quelle croyance, habitude ou comportement a le plus impacté ta vie au cours des cinq dernières années ? \textnormal{\small(1:31:02)}}

C'est l'expérience de la \textbf{maternité} qui a le plus transformé son rapport au monde et sa vision des choses. Loin d'être déconnectée de la crise écologique, elle estime que devenir parent aiguise souvent la conscience de la fragilité et de la valeur de la vie sur Terre, et projette vers un avenir qu'on souhaite accueillant pour ses enfants -- posant la question fondamentale : « Comment aimer nos enfants en contexte d'Anthropocène ? » Elle y rattache plus largement les tâches du \textit{care} -- toutes ces activités où l'on se décentre pour s'intéresser au sort de l'autre -- invisibilisées dans notre société de la performance, mais profondément transformatrices et porteuses d'une autre façon d'habiter le monde. Ces activités confrontent aussi ses réflexions philosophiques à la matière, au sensible et à l'émotion.
